%%%%%%%%%%%%%%%%%%%%%%%%%%%%%%%%%%%%%%%%%%%%%%%%%%%%%%%%%%%%%%%%%%%%%%
% Lớp 8 | Chương I. Đa thức | Phép chia đa thức cho đơn thức
% Dạng 8D1?5-1: Chia đơn thức cho đơn thức
%%%%%%%%%%%%%%%%%%%%%%%%%%%%%%%%%%%%%%%%%%%%%%%%%%%%%%%%%%%%%%%%%%%%%%


% nguồn: De_AI_Toan8_cuoi_hoc_ki_i_1 (nhập 21/11/2026)
\begin{bt}%[8D1V5-1]%[ma:6116cb7f]
	Thực hiện các phép chia đơn thức cho đơn thức sau:
a) $28x^5y^3 : (-7x^2y)$;
b) $(3x^2y)^2 : (xy)^2$.
	\loigiai{
		a) $28x^5y^3 : (-7x^2y) = (28 : (-7)) \cdot (x^5 : x^2) \cdot (y^3 : y) = -4x^3y^2$. \\
b) $(3x^2y)^2 : (xy)^2 = (9x^4y^2) : (x^2y^2) = 9x^2$.
	}
\end{bt}

% nguồn: De_AI_Toan8_cuoi_hoc_ki_i_1 (nhập 21/11/2026)
\begin{bt}%[8D1V5-1]%[ma:7dd97b9e]
	Thực hiện phép tính $A = (-2x^3y^2)^3 : (4x^5y^4)$ rồi tính giá trị của $A$ tại $x = -1$, $y = 3$.
	\loigiai{
		Ta có $A = (-8x^9y^6) : (4x^5y^4) = -2x^4y^2$. \\
Thay $x = -1$, $y = 3$ vào $A$, ta được $A = -2 \cdot (-1)^4 \cdot 3^2 = -2 \cdot 1 \cdot 9 = -18$.
	}
\end{bt}

% nguồn: De_AI_Toan8_cuoi_hoc_ki_i_1 (nhập 21/11/2026)
\begin{bt}%[8D1V5-1]%[ma:86ba8dcd]
	Chứng tỏ rằng hai đơn thức $M = -3x^2y^3$ và $N = 5x^4y^5$ không thể cùng nhận giá trị dương với mọi giá trị của $x$ và $y$.
	\loigiai{
		Xét tích của hai đơn thức: $M \cdot N = (-3x^2y^3)(5x^4y^5) = -15x^6y^8$. \\
Vì $x^6 \ge 0$ và $y^8 \ge 0$ với mọi $x, y$, nên $M \cdot N = -15x^6y^8 \le 0$ với mọi $x, y$. \\
Nếu $M$ và $N$ cùng nhận giá trị dương thì tích $M \cdot N$ phải lớn hơn $0$, điều này mâu thuẫn với $M \cdot N \le 0$. \\
Vậy hai đơn thức $M$ và $N$ không thể cùng có giá trị dương.
	}
\end{bt}
